\section{Renormalization of Non-Local Operators in Lattice Regularization}

The HQET Lagrangian in Eq.~(\ref{Eq:lag}) takes the infinite heavy quark mass limit,
therefore does not contain any mass term. The mass correction to the heavy quark can only receive power divergent contributions~\cite{Maiani:1991az}, which are absent in dimensional regularization.

However, in UV cut-off regularizations such as the lattice regularization,
when going beyond leading-order perturbation theory, the self-energy of the heavy quark introduces a linear divergence which
has to be absorbed into an effective mass counterterm,
\begin{equation}
    \delta {\cal L}_m =  -\delta m\overline{Q} Q
\end{equation}
with $\delta m$ of ${\cal O}(1/a)$ with $a$ the lattice spacing~\cite{Maiani:1991az}. Although lattice regularization breaks Lorentz symmetry and introduces operator mixing, it is not a concern for $O(z_2,z_1)$ because there is no lower-dimension operator that can mix with it in lattice QCD.
In the HQET framework, the appearance of the linear divergence is easy to understand. Given that the heavy quark is infinitely heavy, it behaves like a static color source, whose energy will have a Coulomb-like form $1/r$, and therefore will diverge linearly if the source is a pointlike particle.
In our auxiliary heavy quark Lagrangian, the physical picture as a static color source is lost, whereas the removal of linear divergence can be done essentially in the same way. Since the linear divergence appears only in the self energy of the auxiliary heavy quark, it can be absorbed into the renormalization of the heavy quark mass as in Ref.~\cite{Maiani:1991az}.

Thus, the total heavy-quark Lagrangian now becomes
\begin{equation}\label{Eq:spaceliken}
    {\cal L}_Q = \bar Q (i n\cdot D-\delta m)Q.
\end{equation}
Note that the BRST invariance of the Lagrangian requires a dependence of $\delta m$ on the signature of $n$~\cite{Dorn:1986dt}, which yields a vanishing $\delta m$ for a lightlike $n^\mu$. In the case of a spacelike $n^\mu$, $\delta m$ in Eq.~(\ref{Eq:spaceliken}) is imaginary, we can therefore write $\delta m=i\delta {\bar m}$. After integrating over the heavy fields,  the renormalized operator becomes,
\begin{equation}\label{Eq:cutoffrenorm}
        O_R = Z_{\bar j}^{-1}Z_{j}^{-1}e^{\delta \bar m|z_2-z_1|}\overline{\psi}(z_2) \Gamma L(z_2,z_1)\psi(z_1).
\end{equation}
Thus there is a mass counterterm in the exponent which has a $z$-dependence.
After the mass renormalization, there are at most logarithmic divergences which will be canceled by counterterms from the Lagrangian or from the renormalization of the heavy-light composite operator. The above renormaliztion form has been proposed in a number of papers before~\cite{Musch:2010ka,Ishikawa:2016znu,Chen:2016fxx,Constantinou:2017sej,Alexandrou:2017huk,Chen:2017mzz},
but not proven.

It is worthwhile to point out that the explicit one-loop calculation~\cite{Chen:2016fxx} in lattice regularization indeed shows that the exponential mass renormalization in the form of Eq.~(\ref{Eq:cutoffrenorm}) cancels the linear divergence in the quasi-PDF in Ref.~\cite{Xiong:2013bka} after a Fourier transform to momentum space, where
\beq
\delta \bar m=\frac{2\pi\alpha_s }{3a}.
\eeq

The renormalization presented above works the same when $z_2<z_1$.
